\begin{textsample}{POS Dim 1 – vem\_ed\_18 – Score 25.00 – vem\_ed\_18/t083.txt}  \label{ex:f1_pos_113}
Daniel Nolan, coordenador COIL da University of Minnesota System (EUA), pontuou algumas \textbf{questões} para a \textbf{reflexão} da audiência: Quais \textbf{são} bons modelos para Intercâmbios Virtuais engajados com a comunidade? \textbf{Que} \textbf{habilidades} e recursos \textbf{são} necessários para projetos de Intercâmbios Virtuais \textbf{bem-sucedidos}?
Como engajar professores e alunos nesses \textbf{processos} de parceria?
Em seguida, apresentou dois exemplos de projetos COIL relacionados com a temática do engajamento comunitário.
Osvaldo Succi Junior perguntou como \textbf{começou} esse interesse por \textbf{parte} do professor Dan, \textbf{que} leciona Estudos Germânicos.
Uma dica interessante \textbf{que} ele deu \textbf{foi} “\textbf{começar} \textbf{pequeno} e depois crescer.
O potencial dos parceiros \textbf{que} participam deste evento \textbf{é} enorme”.
Dan Nolan relacionou sua apresentação com a de José Luis Jiménez, sobre a temática dos projetos COIL \textbf{que} se voltam a aspectos humanitários e de ajuda a ONGs \textbf{que} trabalham com refugiados. “Percebemos \textbf{que} os estudantes desenvolvem um senso de responsabilidade com esses projetos”.
E acrescentou: “Humanos \textbf{aprendem} melhor quando há engajamento, porque vivemos em comunidade.
Podemos ajudar as \textbf{pessoas} a \textbf{serem} engajadas e o parceiro internacional \textbf{é} importante nesse \textbf{processo} para \textbf{que} o engajamento \textbf{seja} \textbf{construído} com o passar do tempo”.
Perguntado sobre como desenvolver isso, ele sugere \textbf{que} os professores podem atender aos alunos dentro da instituição para desenvolver Intercâmbios Virtuais. “Minha caixa de ferramentas \textbf{é} a minha Universidade; \textbf{observar} oportunidades já \textbf{desenvolvidas} ajuda muito”.
HUMANISMO E ENGAJAMENTO COMUNITÁRIO José Luis Jiménez, coordenador de Intercâmbios Virtuais da Universidad Católica Andrés Bello (Venezuela), ressaltou a \textbf{importância} da consciência \textbf{crítica} e da aprendizagem transformadora.
A instituição já desenvolveu 95 colaborações com 15 países, envolvendo mais de 2 mil alunos.
O coordenador \textbf{é} também professor de Comunicação e Jornalismo.
Defendeu a abordagem \textbf{interdisciplinar} dos Intercâmbios Virtuais, especialmente nos projetos \textbf{que} buscam soluções para o bem-estar socioambiental, \textbf{utilizando} metodologias \textbf{ativas} como aprendizagem \textbf{baseada} em times (Team-Based Learning) e por projetos (Project-Based Learning) para trazer soluções \textbf{significativas} para a comunidade.
Mencionou a inter-relação entre ativismo, movimentos ligados às \textbf{questões} climáticas e democracia. “Exploramos o potencial dos projetos COIL para \textbf{promover} os Objetivos de Desenvolvimento \textbf{Sustentável} (ODS) da ONU em países com \textbf{crises} humanitárias – como \textbf{todos} \textbf{sabem}, esse \textbf{é} o caso da Venezuela”, comentou.
No \textbf{momento}, a Universidad Católica Andrés Bello está realizando um PCI com a Fatec Barueri sobre a \textbf{questão} dos Yanomami, \textbf{que} tiveram seu território dividido pelas fronteiras do Brasil e da Venezuela e passam por uma situação \textbf{crítica} do \textbf{ponto} de vista humanitário, como se vê na cobertura de imprensa.

% matched lemmas: aprender, ativo, basear, bem-sucedido, começar, construir, crise, crítico, desenvolvir, habilidade, importância, interdisciplinar, momento, observar, parte, pequeno, pessoa, ponto, processo, promover, que, questão, reflexão, saber, ser, significativo, sustentável, todo, utilizar
\end{textsample}
